% ==============================================================================
% MẪU TẠO TABLE MỚI TRONG TIKZ ERD
% ==============================================================================
% Bạn có thể copy khối mẫu này vào file brewlite-erd.tex hoặc file .tex riêng

% 1. Khai báo Node Table:
% \node[table, right=3.5cm of <node_goc>] (<ten_node>) {
%     \textbf{\textcolor{tableheadertext}{<TenBang>}}
%     \nodepart{two}
%     \underline{\textcolor{pkcolor}{id}} : Int \textbf{[PK]} \\
%     \textcolor{fkcolor}{foreignId} : Int \textbf{[FK]} \\
%     tenThuocTinh : String \\
%     trangThai : Int (def 0) \\
%     createdAt : DateTime
% };

% 2. Khai báo Quan hệ (Relationship) giữa 2 bảng:
% - Đường thẳng nối bảng cha (1) sang bảng con (N):
% \draw[relation] (<bang_cha>.east) -- 
%     node[cardinality, above, pos=0.08] {1} 
%     node[cardinality, above, pos=0.92] {N} 
%     (<bang_con>.west);

% - Đường cong nếu các bảng nằm lệch hàng / tránh chồng chéo:
% \draw[relation] (<bang_cha>.south) to[out=-90, in=180] 
%     node[cardinality, right, pos=0.08] {1} 
%     node[cardinality, above, pos=0.92] {N} 
%     (<bang_con>.west);
